\section{问题定义与设计目标}

\subsection{BLoP 基线与扩展目标}

本文使用 BLoP 作为基线可信审计模型~\cite{Zhang2025BloPAT}。BLoP 通过可信计算、可信代理、事件监测和 PBFT 风格确认提供完整信任链。其自然运行模式是全节点确认：所有符合条件的可信节点共同确认关键审计事件。当高风险数据使用、强问责要求和低频事件占主导时，这是一种合适的默认模式。

扩展目标则不同。在数据流通中，数据连接器、存储节点、计算任务和监测代理可能持续产生常规审计事件。如果所有常规事件都使用同一全节点路径处理，链上记录和确认消息可能比单个事件本身的价值增长得更快。因此，本文研究问题为：

\begin{quote}
一个 BLoP 风格的可信审计系统，能否在保留高风险或争议事件全节点确认的同时，为常规事件支持轻量化、风险感知的处理路径？
\end{quote}

\subsection{系统问题}

考虑一个由多个组织参与的联盟式数据流通场景。数据拥有方注册数据对象，将其拆分为分片，并将分片或计算证据存储在链下节点。审计者周期性地挑战被选中的分片或证据对象。联盟链记录承诺、审计摘要、恢复事件和信任分更新。系统应在不暴露明文数据、不向链上提交大体量证据的前提下，校验数据或证据完整性。

TrustAuditFlow 解决的问题是设计一种低链上成本的校验机制，使其能够检测损坏或缺失的链下证据，在可能时触发恢复，记录责任，并更新未来节点选择。该机制还必须明确何时轻量委员会确认足够，何时需要全节点确认。

\subsection{威胁模型}

本文假设许可环境中参与组织和节点均已注册，但部分节点可能故障或拜占庭。链下存储节点可能删除分片、返回过期数据、延迟响应或提交不一致证据。审计节点可能遗漏计划审计、错误报告审计状态，或与存储节点串谋。联盟链在 PBFT 风格假设下容忍至多 $f$ 个拜占庭节点。委员会确认仅在可用投票节点满足所需法定人数时被接受。

该设计不处理密码学原语被破坏、委员会被自适应攻击者完全控制、针对 TEE 的硬件侧信道攻击，或上游数据生成过程中的业务语义欺诈。这些边界是刻意设置的：本文关注审计流程和证据治理层，而不声称端到端业务语义正确性。

\subsection{设计目标}

TrustAuditFlow 围绕六个目标设计。

\begin{itemize}
  \item \textbf{低链上成本。} 链上存储承诺、Merkle 根、审计状态、恢复状态和信任更新，而非原始数据或大型证明。
  \item \textbf{过程可追踪。} 审计请求、异常检测、恢复动作和节点惩罚必须锚定在防篡改时间线中。
  \item \textbf{恢复闭环。} 当损坏分片数量处于纠删码恢复阈值以内时，审计失败应触发恢复和承诺更新，而不只是返回失败。
  \item \textbf{轻量委员会确认。} 信任分驱动的委员会作为 BLoP 可信节点集合之上的可选路径，用于确认常规事件。
  \item \textbf{全节点回退。} 高风险对象、委员会信任不足或争议审计事件回退到 BLoP 风格全节点确认。
  \item \textbf{明确证据边界。} TrustAuditFlow 不替代 BLoP 的可信计算层或 PDP/PoR 密码学协议，而是组织这些组件如何被锚定、确认并连接到恢复过程。
\end{itemize}

